% Cote 119: offprints, reports, courses, CV and letters (typed and printed part): archive cote 119, site page 177 (PDF page 178).
% typescript. Draft compared: none.
% AI-transcribed by Claude Opus 5.5 (claude-opus-5-5, adaptive thinking, effort high), Message Batches API, batch r119_full (msgbatch_01DwZNm5Y7YsKBWyAN4sFs8w), 2026-11-13; not yet checked by eye.
\scanpage{119/177}
\mspage{- 16 -}

6.2. \uline{Définition}.

On dit que deux chemins
\[ \gamma_0 : I \to U \quad \text{et} \quad \gamma_1 : I \to U \qquad \text{(U ouvert)} \]
sont homotopes dans U s'il existe une application continue
\[ \delta : I \times I \to U \qquad (t,u) \to \delta(t,u) \]
telle que \qquad $\delta(t,0) = \gamma_0(t)$

(6.2.1)

\qquad\qquad $\delta(t,1) = \gamma_1(t)$

Si $\gamma_0$ et $\gamma_1$ ont même origine et même extrémité on dit que $\gamma_1$ et $\gamma_1$
sont homotopes avec extrémités fixes

(6.2.2.) \qquad si de plus : \qquad $\delta(0,u) = \gamma_0(0) = \gamma_1(0)$

\qquad\qquad $\delta(1,u) = \gamma_0(1) = \gamma_1(1)$

Si $\gamma_0$ et $\gamma_1$ sont des chemins fermés\unsure{,} on dit qu'ils sont homotopes comme
chemins fermés si l'homotopie $\delta$ vérifie.
\[ \delta(t,0) = \gamma_0(t) \]
(6.2.1.)
\[ \delta(t,1) = \gamma_1(t) \]

(6.2.3.) \qquad $\delta(0,u) = \delta(\unsure{1},u)$

autrement dit s'ils sont homotopes et si pour tout u : $t \to \delta(t,u)$ est un
chemin fermé. (Si $\gamma_1(t)$ est constante on dira que le chemin fermé $\gamma_0(t)$
est homotope à un point).

6.3. \uline{Définition}.

Un ouvert U est simplement connexe si tout chemin fermé dans U
est homotope à un point (comme chemin fermé bien sûr).et si U est connexe.

6.4. \uline{Exemples} : Un disque est simplement connexe. Un ensemble convexe et plus
généralement un ensemble étoilé est simplement connexe. (Un ensemble U est
étoilé par rapport à $Z_0 \in U$ siopour tout $Z \in U$ le segment $[Z_0,Z]$ est
contenu dans U). Il suffit de considérer l'homotopie
\[ \delta(t,u) = u\,(\gamma(t) - Z_0) + Z_0 \ . \]
Le plan complexe privé des points réels négatifs est donc simplement connexe.
C - \{0\} ne l'est pas. Une couronne n'est pas simplement connexe.

\mspage{177}
